% Physical unit 043; accepted Session 03 after fresh 400-dpi authority replay.
\GordanUnitStart{043}{printed page 38}
\GordanTopFolio{38}
\begingroup\small
Hence:
\[
\left\{\tg^2l_0''+\tg l_1\tg l_2\right\}\left\{\pi-\psi-e^2{\pi\over2}\cos l_0''\right\}^2=\left\{\tg l_1+\tg l_2\right\}^2
\]
\[
 \cos^2l_1\cos^2l_2\left\{{1\over\cos^2l_0''}-{\cos(l_1+l_2)\over\cos l_1\cos l_2}\right\}\left\{\pi-\psi-{e^2\over2}\pi\cos l_0''\right\}^2=\sin^2(l_1+l_2)
\]
\[
 \cos l_1\cos l_2\left\{{\cos l_1\cos l_2\over\cos^2l_0''}-\cos(l_1+l_2)\right\}\left\{\pi-\psi-{e^2\pi\over2}\cos l_0''\right\}^2=\sin^2(l_1+l_2)
\]
or, putting
\[
\text{in place of }\cos l_2:\ \cos l_1;\quad \text{in place of }\cos(l_1+l_2):\ 1;\quad \text{in place of }\sin(l_1+l_2):\ l_1+l_2,
\]
\[
 \cos^2l_1\left\{{\cos^2l_1\over\cos^2l_0''}-1\right\}\left\{\pi-\psi-e^2{\pi\over2}\cos l_0''\right\}^2=(l_1+l_2)^2
\]
\[
 \cos l_1\cot T\left\{\pi-\psi-{e^2\over2}\pi\cos l_0''\right\}=l_1+l_2.
\]
In this formula $T$ is known to signify the azimuth of the geodetic line $l_0''$ at the point $A$; therefore in $l_0''$ the formula holds
\[
 \sin T\cos l_1=\cos l_0''.
\]
Putting this value of $\cos l_0''$ in our equation, there results the formula
\[
 \cos l_1\cot T\{\pi-\psi\}-e^2{\pi\over2}\cos^2l_1\cos T=l_1+l_2,
\]
from which $T$ and likewise $l_0''$ can easily be computed.

That equation is sought from these geometrical considerations. The geodetic line $l_2$, produced beyond the point $B$ until it cuts the circle \gq{$-l_0$} (for example in $C$), with the circle $-l_1$ and with the meridian of the point $B$ (which curves we make cut in $D$) forms a right triangle whose sides are so short that the triangle can be understood and taken as plane:
\[
 {BD\over CD}=\tg BCD \qquad\text{cf. chap. I \S\ 9.}
\]
But, as is easily understood,
\[
 BD=a\int_{-l_1}^{l_2} dl\sqrt{1-e^2\cos^2l}
\]
\[
 CD=a\cos l_1\{2\mathcal A(l_0)-\psi\}
\]
and the angle $BCD$ makes the complement of $T$ itself. Substituting these into our formula, we obtain
\[
 \cot T\cos l_1\{2\mathcal A(l_0)-\psi\}=\int_{-l_1}^{l_2}dl\sqrt{1-e^2\cos^2l}
\]
or, neglecting the eccentricity under the sign of integration, and, in the series expansion of the period longitude in Chap. III, taking account only of the first two terms,
\[
 \cot T\cos l_1\left\{\pi-e^2{\pi\over2}\cos l_0-\psi\right\}=l_1+l_2,
\]
which was to be demonstrated.
\endgroup
